% =============================================================================
% CARTOGRASS 2026 poster -- TEXTS: context row and concept row.
% -----------------------------------------------------------------------------
% Row 2: two motivation panels flanking two large images (Why / Two gaps).
% Row 3: concept panel, OBJECTIVE box (one question, large type) and the small
%        pilot-result box with a QR.
% Use: \begin{pblock}{\cgWhyTitle}\cgWhyBody\end{pblock}, and so on.
% Word counts: words with a letter as pdftotext reads them, titles excluded.
% Budget for the whole sheet (design review, 2026-10-16): about 550 words in
% all, captions, footer and the text inside figures included; each motivation
% panel about 30 words.
% Audience pass (2026-10-19): plain words for seagrass ecologists, cartographers
% and managers; the words of the method flow chart ("delay", "out of the
% water", "share of time"); no statistics or computing terms left unexplained.
% Citation numbers follow first appearance on the poster; see text_footer.tex.
% =============================================================================
\ifdefined\cgContextTextLoaded\endinput\fi
\def\cgContextTextLoaded{}

% --- Motivation panel 1 ------------------------------------------- 31 words
% Definition as in the submitted abstract ("spends exposed"), in the flow
% chart's words (card 5: "share of time ... out of the water"); the last
% sentence is the bridge to "Two gaps" on the other side of the images.
% Round 3 (2026-10-19): "(emersion time)", the seagrass ecologists' own term
% (also "exposure time"), so "emerged time" does not read as a new quantity.
% TODO(authors): the claim "a main driver" carries no reference; a supporting
% one (chosen by the ecologist co-author) needs about 8 mm the references
% block does not have (48 of 52 mm used): add it if another is shortened.
\newcommand{\cgWhyTitle}{Why emerged time?}
\newcommand{\cgWhyBody}{%
  A main driver of where \emph{Zostera noltei} grows is \term{emerged time}
  (emersion time): the share of time the ground spends out of the water, set
  by its elevation and the local tide.}
% The photo slot under it (\cgphoto in blk_ctx_left.tex): the label says
% which photo is wanted until figures/cg/photo_zostera.(jpg|png) exists.
\newcommand{\cgWhyPhoto}{\emph{Zostera noltei} exposed at low tide,
  Ría de Villaviciosa (authors' own)}

% --- Motivation panel 2 ------------------------------------------- 29 words
% "About 2 hours": the largest delays measured inside the sites, behind the
% tide at the mouth. Recomputed for 2025 against Borkum (science review,
% 2026-10-16): Nieuwe Statenzijl (Dollard) best-aligning shift +96 min, its
% low waters +124 to +133 min; Pogum high waters +96 to +102 min; Harlingen
% (Wadden Sea) best constant lag +93 min (+88 min in the brief's gauge
% table). "Up to 1.5 hours" understated the maximum. The concept figure (b1)
% shows the Ems at +58 / +86 min for one high water, consistent with "up to".
% "Elevation map", not "bathymetry" (2026-10-19): the flats are land at low
% tide, and the poster's title says topography.
% Round 2 (2026-10-19): global models such as EOT20 are fitted to satellite
% altimetry, so "calibrated at the coast" was wrong; what fails inside is
% resolution (they give the tide at the open coast). "1.5-2 hours": the
% whole-curve delays MAREA corrects reach about 1.5 h at the gauges
% (Nieuwe Statenzijl +96 min, Harlingen +93 min); 2 h is reached only by the
% Dollard's low waters (+124 to +133 min).
\newcommand{\cgGapsTitle}{Two gaps}
% Its punch line (\cgquote), the abstract's "To close both gaps, we present
% MAREA" as intent, not result: "addresses", not "closes".
\newcommand{\cgGapsQuote}{MAREA addresses both with the Sentinel-2
  archive.}
% Round 3 (2026-10-19): "tidal flats", not "wetlands", which brought in
% marshes (Zostera noltei lives on the flats), and "up-to-date", not
% "detailed": PNOA LiDAR covers all of Spain and the Vaklodingen the whole
% Dutch coast (both are checks at four of the five sites), but lidar is flown
% at whatever tide stage and depth surveys come every few years, so what is
% missing is a current map of the flats. "Tidal", not "intertidal": the
% longer word left "map." alone on a third line.
% Round 4 (2026-10-19): (a) "estuaries and tidal basins", not "rias and
% estuaries": the 1.5-2 h hold for the Dollard and the Wadden basin
% (Harlingen +93 min), while the poster's own rias show little delay
% (Villaviciosa 0-21 min, Ferrol short). (b) The consequence is said, the
% mechanism only (no size of any gain, which stays with Results): a map made
% with the global models' coastal tide misjudges when the flats flood ("made
% with them": the models; "with the coastal tide" took a sixth line, 1.6 mm
% past the panel); the concept panel's pull quote ("If the tide arrives late,
% the pixel floods late") picks it up.
\newcommand{\cgGapsBody}{%
  \term{Unknown elevation.} Most tidal flats lack an up-to-date elevation
  map.\par
  \term{A late tide.} Global tide models give the tide at the open coast;
  inside estuaries and tidal basins it arrives up to \emphnum{1.5--2~hours}
  later, so maps made with them misjudge when flats flood.}

% --- Concept panel ------------------------------------------------ 29 words
% [1, 2] = Bishop-Taylor et al. 2019; Granadeiro et al. 2021: the wet/dry
% versus tide principle. Never claim a first: [2] already estimates lags.
% "Many times": MAREA's 2023-2025 records hold 54-106 scenes per site and the
% pixel figure 95 looks (Villaviciosa, 97 scenes); "hundreds" holds only for
% 2016-2025.
% "A yes-or-no tide gauge" (2026-10-19), not "one-bit": no computing terms.
% Never claim a first (round 4): the pull quote below credits [2] for the
% delay itself.
% The title breaks by hand after "satellite" (36 pt, 165 mm measure).
\newcommand{\cgSeeTitle}{What does a satellite\\see of the tide?}
% Round 3: "Mostly wet or dry", not "Only": the elevation step fits a curve
% to the continuous water index (pyintertidal.marea.invert_series; the flow
% chart's card 4 and "A pixel is a tide gauge"); only the delay fit uses the
% looks as wet or dry.
\newcommand{\cgSeeBody}{%
  Mostly wet or dry: a pixel is wet when the water stands above its
  ground~[1,\,2]. Seen many times, at different tide stages, each pixel
  becomes a \term{yes-or-no tide gauge}.}
% Its punch line, set large under it (\cgquote): the bridge to the figure.
% Round 4 (2026-10-19): credited to [2]. Granadeiro et al. 2021 is titled
% "Using Sentinel-2 images to estimate topography, tidal-stage lags and
% exposure periods": the same three outputs as MAREA, and the project's own
% docs/prior_art_granadeiro.md says [2] showed that the Sentinel-2 archive
% holds the delays. [2] fits one delay field in lon/lat from flood and ebb
% scenes on a sample of mid-elevation pixels; MAREA fits one delay per 3-km
% section along the channel from every wet or dry look (the OBJECTIVE names
% that difference: "section by section along the channel").
\newcommand{\cgSeeQuote}{If the tide arrives late, the pixel floods
  late~[2].}

% --- Objective: one question in large type ------------------------ 25 words
\newcommand{\cgObjectiveLabel}{Objective}
% Round 4 (2026-10-19): "section by section along the channel" names what
% MAREA adds to [2] (one delay per 3-km section, zero at the mouth, never
% decreasing inland), and drops "inside rias and estuaries" (the rias show
% little delay; the Wadden site is a tidal basin): "along the channel" says
% where the delay is measured.
\newcommand{\cgObjective}{Can the Sentinel-2 archive measure, section by
  section along the channel, how late the tide arrives and, with that delay
  corrected, map intertidal elevation and emerged time?}
% Why a few decimeters matter (2026-10-19), set under the question in smaller
% type (blk_con_objective.tex): the number of \cgCapHours (b2, Gijon tide
% gauge, 2024: -0.15 to +0.15 m around mean sea level = +1.02 h in the air per
% tide). Tide-gauge data, not a MAREA result. ------------------- 28 words
% Round 4 (2026-10-19): "a flat 30 cm higher spends", the ecologist's
% picture (a flat higher up the shore) in the words of the sheet's own
% definition ("the share of time the ground spends out of the water").
\newcommand{\cgObjectiveWhy}{\textbf{Why decimeters matter:} around mean
  sea level, a flat 30\,cm higher spends about 1\,hour more out of the
  water per tide (Gijón tide gauge, 2024).}

% --- Pilot result (previous work), small box with a QR ------------ 25 words
% 194 points and 0.19 m are from the submitted abstract (which also says "5%
% of the spring range": see round 4 below). 194 is the number of RTK points SCORED (on mapped pixels), not the
% size of the survey (361 points; 270 development points), so the count is
% given here and not in the materials list.
% "Typical elevation error 0.19 m (offset removed)" (2026-10-19): the spread
% of the differences once their offset is removed (centered RMSE; the
% validation centres on the median), in Results B's words ("typical error");
% "precise GPS points" is the materials list's gloss of RTK-GNSS, so the
% first number on the sheet needs no term explained later. The box holds
% five lines (151 mm at 27.9 pt): one more overflows it.
% The delay barely matters at Villaviciosa (0-21 min; centered RMSE 0.191 m
% without it, 0.193 m with it, per the science review of 2026-10-16): the box
% says so.
% TODO(authors): when Results B is filled from validation v2, print one
% elevation number on the whole sheet, or keep this box as the pilot.
% TODO: decide the QR target (validation summary? abstract page?).
% Round 4 (2026-10-19): "6% of the mean spring range" (5.6 % rounded; the
% box holds five lines, so no "about"), not the abstract's "5% of the spring
% tidal range". The cached EOT20 series at
% Villaviciosa (tools/cg/cache/villaviciosa_eot20_hourly_2023.csv, harmonic
% fit) gives M2 = 1.25 m and S2 = 0.45 m, so the mean spring range
% 2 (M2 + S2) is 3.41 m and 0.19 m is 5.6 % of it; 5 % holds only against a
% larger range (the 90th percentile of the daily ranges, 3.74 m, gives 5.1 %).
\newcommand{\cgProofTitle}{Pilot result}
\newcommand{\cgProofBody}{%
  Ría de Villaviciosa, 194 precise GPS points: typical elevation error
  \emphnum{0.19\,m} (offset removed), 6\% of the mean spring range. Little
  delay there: a test of the elevation step.}
\newcommand{\cgProofQRLabel}{Validation details}
